\documentclass[10pt,twocolumn]{article}
\usepackage[letterpaper,margin=0.58in,columnsep=0.24in]{geometry}
\usepackage[T1]{fontenc}
\usepackage{lmodern}
\usepackage{microtype}
\usepackage{enumitem}
\usepackage{xcolor}
\usepackage{hyperref}
\usepackage{listings}
\usepackage{titlesec}
\usepackage{fancyhdr}
\usepackage{flushend}

\hypersetup{
  colorlinks=true,
  linkcolor=blue!50!black,
  urlcolor=blue!50!black,
  pdftitle={DeepStack MICRO 2026 Artifact Appendix},
  pdfauthor={DeepStack authors}
}
\setlength{\parindent}{0pt}
\setlength{\parskip}{2.2pt}
\setlist{nosep,leftmargin=*,topsep=1.5pt}
\titlespacing*{\section}{0pt}{5pt}{1.5pt}
\titlespacing*{\subsection}{0pt}{3.5pt}{1pt}
\titlespacing*{\subsubsection}{0pt}{2.5pt}{1pt}
\titleformat{\section}{\bfseries\normalsize}{\thesection}{0.4em}{}
\titleformat{\subsection}{\bfseries\small}{\thesubsection}{0.4em}{}
\titleformat{\subsubsection}{\bfseries\itshape\small}{\thesubsubsection}{0.4em}{}
\pagestyle{fancy}
\fancyhf{}
\lhead{DeepStack Artifact Appendix}
\rhead{DeepStack}
\cfoot{\thepage}
\renewcommand{\headrulewidth}{0.25pt}
\lstset{
  basicstyle=\ttfamily\scriptsize,
  frame=single,
  framerule=0.3pt,
  breaklines=true,
  columns=fullflexible,
  aboveskip=3pt,
  belowskip=3pt,
  xleftmargin=1pt,
  xrightmargin=1pt,
  backgroundcolor=\color{black!2}
}
\newcommand{\code}[1]{\texttt{#1}}

\begin{document}
\appendix
\section{Artifact Appendix}

\subsection{Abstract}

This artifact provides the implementation and automated evaluation workflow
for DeepStack's analytical modeling and design-space exploration (DSE) of
distributed 3D DRAM-stacked LLM accelerators. It reruns the fixed
configurations underlying paper Figs.~8--12, Figs.~13--21, and Table~4; checks
them against bundled, checksummed GPU and simulator references; and
regenerates the paper-facing plots. No GPU, model weights, or external
simulator is required.

\subsection{Artifact check-list (meta-information)}

{\small
\begin{itemize}
  \item {\bf Algorithm:} analytical latency, throughput, energy, thermal, and
        network modeling with hierarchical DSE.
  \item {\bf Program:} Python~3.11 package and staged \code{run.sh} workflow.
  \item {\bf Compilation/Binary:} no local compilation; four licensed,
        prebuilt CPython~3.11 x86-64 extensions are included.
  \item {\bf Model:} DeepSeek-V3, Qwen3-235B-A22B, Llama-70B, and
        kernel-level workload descriptions; no model weights.
  \item {\bf Data set:} checksummed GPU/vLLM and ASTRA-Sim/NS-3 references,
        fixed-configuration manifests, routing traces, and a recorded
        301{,}728-row decode-DSE population.
  \item {\bf Run-time environment:} x86-64 Linux, GNU Bash, Conda,
        CPython~3.11, and glibc~2.29 or newer.
  \item {\bf Hardware:} 32 CPU cores and 64~GiB RAM recommended; no GPU,
        FPGA, or accelerator.
  \item {\bf Metrics:} weighted MAPE, scaled tokens/s, bandwidth, speedup,
        temperature, and area feasibility.
  \item {\bf Output:} validated CSVs, run/verification manifests, and 16
        figure groups in PNG and PDF.
  \item {\bf Experiments:} 15 fail-fast stages plus independent verification.
  \item {\bf Disk space:} approximately 5~GiB including the environment.
  \item {\bf Preparation time:} approximately 5--10 minutes.
  \item {\bf Experiment time:} 12.2 minutes on the 64-core/128-thread
        packaging host with 32 workers; 30--45 minutes estimated on 32 cores.
  \item {\bf Publicly available?:} GitHub and archived DOI.
  \item {\bf Licenses:} Apache-2.0 for source and data; bundled extensions use
        LicenseRef-DeepStack-AE-Binary-1.0.
  \item {\bf Automation:} repository-local scripts; no external framework.
  \item {\bf Archived DOI:}
        \href{https://doi.org/10.5281/zenodo.21351389}{10.5281/zenodo.21351389}.
\end{itemize}
}

\subsection{Description}

\subsubsection{How to access}

Clone \href{https://github.com/tile-ai/DeepStack}{github.com/tile-ai/DeepStack}
(branch \code{ae}) or download the DOI archive above. The root
\code{README.md} is the reviewer guide.

\subsubsection{Hardware dependencies}

Use a 64-bit x86 Linux host with 32 cores, 64~GiB RAM, and 5~GiB free
storage. Lower core counts work with fewer workers and longer runtime. GPU and
simulator measurements are bundled reference inputs and are not rerun.

\subsubsection{Software dependencies}

The artifact requires GNU Bash, Conda, CPython~3.11, and glibc~2.29 or newer.
\code{./setup.sh} installs pinned dependencies. Four licensed prebuilt
extensions provide protected model components; no compiler, CUDA, external
simulator, or proprietary EDA tool is needed.

\subsubsection{Data sets and models}

All required references, configurations, routing traces, and plotting data are
included and covered by \code{SHA256SUMS}. The artifact contains workload
descriptions and routing metadata, but no prompts, weights, or hidden states.
The complete area model and die-area breakdown are not distributed. Reference
feasibility exposes only a feasible-SM count from a finite compiled capacity
table.

\subsection{Installation}

\begin{lstlisting}[language=bash]
./setup.sh
conda run --no-capture-output -n deepstack-ae \
  ./run.sh doctor
\end{lstlisting}

Installation succeeds when the second command prints \code{doctor: PASS}.

\subsection{Experiment workflow}

\begin{lstlisting}[language=bash]
conda run --no-capture-output -n deepstack-ae \
  ./run.sh quick --model-version both
conda run --no-capture-output -n deepstack-ae \
  ./run.sh reproduce --workers 32 \
  --model-version both
./verify.sh results/reproduce
\end{lstlisting}

Checked-in renderings under \code{prebuilt/figures/} are inspection-only
previews. The workflow creates an independent result bundle under
\code{results/reproduce/}; neither the model nor verifier reads the prebuilt
directory.

Regenerated paper-facing plots use \code{paper\_legacy};
\code{current\_corrected} is separate diagnostic output. Paper Fig.~9
additionally audits the frozen paper statistic separately from its legacy
rerun. No reported statistic mixes the two paths.

\subsection{Evaluation and expected results}

A successful full run prints \code{workflow: PASS} and
\code{Overall verification: PASS}, then produces 57 validated CSVs, 16 status
markers, and 16 figure groups. Artifact stage IDs retain the accepted-revision
numbering: paper Figs.~11, 12, and 13--21 map to \code{fig12}, \code{fig13},
and \code{fig15}--\code{fig23}; the first three validation stages are
unchanged.

\begin{itemize}
  \item {\bf Validation (paper Figs.~8--12):} H100 weighted MAPE
        $8.366\%$; B200 frozen-paper MAPE $12.918976\%$ and legacy-rerun MAPE
        $12.724500\%$; MI325X mean per-model weighted MAPE $13.095\%$;
        sparse-attention legacy weighted MAPE $2.612\%$; and verified NoC
        closure including the reported $108{,}000\times$ endpoint ratio.
  \item {\bf Co-design (paper Figs.~13--21 and Table~4):} the workflow
        regenerates the selected winners and sweeps, including the nine-layer
        bandwidth peak, $39.62\%$ decrease by layer 12, and Table~4 speedups
        of $2.7906\times$ and $9.4744\times$.
\end{itemize}

Exact checks and tolerances are in each stage's \code{summary.csv} and
\code{docs/result\_matrix.md}. The full $\sim$$2.5\times10^{14}$-point search
and the non-distributable hardware-emulation reference are outside the
reproduced scope.

\subsection{Experiment customization}

Custom architecture, energy, NoC, and DSE interfaces and runnable synthetic
examples are documented in \code{README.md} and \code{tests/}. They include
caller-owned capacity providers or area estimators.

\subsection{Notes}

On failure, inspect the stage log named in \code{run\_manifest.csv}. For
memory pressure, rerun with \code{--workers 16}.

\end{document}
